\documentclass[UTF8,11pt,a4paper,fontset=fandol]{ctexart}
\usepackage[margin=20mm]{geometry}
\usepackage{amsmath,amssymb,booktabs,array,tabularx,enumitem,setspace,xcolor,hyperref,xurl,fancyhdr}
\setstretch{1.08}
\setlength{\emergencystretch}{3em}
\setlength{\headheight}{15pt}
\setlength{\tabcolsep}{5pt}
\renewcommand{\arraystretch}{1.2}
\setlist{nosep,leftmargin=2em}
\hypersetup{colorlinks=true,linkcolor=blue!45!black,urlcolor=blue!45!black}
\providecommand{\WorkloadName}{工作负载}
\pagestyle{fancy}\fancyhf{}
\fancyhead[L]{\small Table II(b)：算子级 RI}
\fancyhead[R]{\small\WorkloadName}
\fancyfoot[C]{\thepage}
\newcommand{\qs}{Q_{\mathrm S}}
\newcommand{\qr}{Q_{\mathrm R}}
\newcommand{\RI}{\mathrm{RI}}
\newcommand{\para}[1]{\par\smallskip\noindent\textbf{#1}\quad}
\newcommand{\ModelTableSetup}{\renewcommand{\arraystretch}{1.2}\setlength{\tabcolsep}{4pt}}
\newcommand{\ResultTableSetup}{\renewcommand{\arraystretch}{1.45}\setlength{\tabcolsep}{4pt}}
\newcommand{\DisplayNote}{\par\smallskip{\noindent\footnotesize $1\mathrm{K}=1024$，$1\mathrm{M}=1024^2$。RI 至少为 1 时保留一位小数，小于 1 时保留三位有效数字；精确值见机器数据。\par}\smallskip}
\author{}\date{2026 年 10 月 2 日}

% Compact, identical title block keeps the mathematical introduction on one page.
\makeatletter
\renewcommand{\maketitle}{\thispagestyle{plain}\begin{center}
{\LARGE\bfseries\@title\par}\vspace{8pt}{\normalsize\@date\par}
\end{center}\vspace{3pt}}
\makeatother

\renewcommand{\WorkloadName}{FFN／MoE}
\title{\bfseries\WorkloadName 的输入复用与 RI}
\begin{document}
\maketitle

\section{对象与公式}
每个窗口取一个 Dense FFN 或一个路由专家，完整装载其 gate、up、down
三矩阵一次，随后累计服务 $U$ 个向量。$U$ 为复用次数（reuse count），包含首次使用；可跨多个 batch 或请求累计。
沿用 Table II(a) 的等宽口径，所有计数角色均为每元素 1 Byte。

设模型隐藏维为 $D$，FFN／单专家中间宽度为 $F$。三矩阵及累计输入的关系为
\[
 \begin{gathered}
 W_g,W_u\in\mathbb{R}^{F\times D},\qquad W_d\in\mathbb{R}^{D\times F},\qquad
 X\in\mathbb{R}^{U\times D},\\
 Z=\operatorname{SiLU}(XW_g^{\mathsf T})\odot(XW_u^{\mathsf T}),\qquad
 Y=ZW_d^{\mathsf T}.
 \end{gathered}
\]
在矩阵阶段入口计数：gate/up 共用同一源输入 $X$，只计 $UD$ Byte；
$Z$ 是 down 的新输入，另计 $UF$ Byte。
因此，三个分项各自的输入为 $UD,UD,UF$，汇总时扣除一次 $UD$ 的共享重叠。
\[
 \boxed{\qs=U(D+F),\qquad \qr=3DF,\qquad
 \RI=\frac{U(D+F)}{3DF}.}
\]
$\qs$ 与 $\qr$ 的单位均为 Byte，RI 为二者之比。
初始有效权重容量为 0，末态保留三矩阵，容量为 $3DF$ Byte。
当前输出不直接加进 $\qs$；若它成为下一矩阵阶段的输入，则在该阶段计数。
这里的 gate 是 SwiGLU 分支，路由矩阵不在此对象内；
单专家结果不乘 top-k、专家总数、共享专家数或层数。

\section{模型信息}
下表保留固定 checkpoint 的原生 $D,F$，层号从 0 起。
前两项取一个 Dense FFN，后四项取所选 MoE 层中的一个路由专家。
\begin{center}\small\ModelTableSetup
\begin{tabularx}{\linewidth}{@{}p{0.32\linewidth}*{4}{>{\centering\arraybackslash}X}@{}}
\toprule
模型 & 所选对象 & 层号 & $D$ & $F$ \\
\midrule
Qwen3.5-2B & Dense FFN & 3 & 2048 & 6144 \\
Ministral 3 8B (2512) & Dense FFN & 0 & 4096 & 14336 \\
Qwen3.6-35B-A3B & 单路由专家 & 3 & 2048 & 512 \\
Tencent Hy3 (295B) & 单路由专家 & 1 & 4096 & 1536 \\
Ling-1T & 单路由专家 & 4 & 8192 & 2048 \\
MiMo-V2.5-Pro & 单路由专家 & 7 & 6144 & 2048 \\
\bottomrule
\end{tabularx}
\end{center}

\newpage
\section{代入结果}
六模型各扫描 $U\in\{1,16,128,1\mathrm{K},16\mathrm{K}\}$，得到 30 个有限工况。
表内为小数 RI；三矩阵的一次写入量在同一模型的各列中保持不变。
\begin{center}\small\ResultTableSetup
\begin{tabularx}{\linewidth}{@{}p{0.32\linewidth}*{5}{>{\centering\arraybackslash}X}@{}}
\toprule
模型 & $U=1$ & $U=16$ & $U=128$ & $U=1\mathrm{K}$ & $U=16\mathrm{K}$ \\
\midrule
Qwen3.5-2B & $0.000217$ & $0.00347$ & $0.0278$ & $0.222$ & $3.6$ \\
Ministral 3 8B (2512) & $0.000105$ & $0.00167$ & $0.0134$ & $0.107$ & $1.7$ \\
Qwen3.6-35B-A3B & $0.000814$ & $0.0130$ & $0.104$ & $0.833$ & $13.3$ \\
Tencent Hy3 (295B) & $0.000298$ & $0.00477$ & $0.0382$ & $0.306$ & $4.9$ \\
Ling-1T & $0.000203$ & $0.00326$ & $0.0260$ & $0.208$ & $3.3$ \\
MiMo-V2.5-Pro & $0.000217$ & $0.00347$ & $0.0278$ & $0.222$ & $3.6$ \\
\bottomrule
\end{tabularx}
\end{center}
\DisplayNote

例如，Qwen3.5-2B 在 $U=1\mathrm{K}$ 时，输入量为 $8\,388\,608$ Byte，
一次写入量为 $37\,748\,736$ Byte，故
\[
 \RI=\frac{1024(2048+6144)}{3\times2048\times6144}\approx0.222.
\]

\section{结果说明}
固定 $D,F$ 后，$\qr$ 与末态有效权重容量固定，$\qs$ 和 RI 随 $U$ 线性增长。
这与 QKV 都是一次装载后的复用关系。专家的选择与替换策略可能使其实际 $U$ 较小；
替换更频繁时，同一装载服务的向量数更少。本表扫描 $U$，不从模型尺寸推断实际替换频率。

固定 $U$ 时，$\RI=(U/3)(1/D+1/F)$。
本表中 Qwen3.6 的单专家宽度 $F=512$，对应最高 RI；
Ministral 的 $D,F$ 组合对应最低 RI。
比较依据是所选三矩阵的维度，模型名中的总参数标签不进入公式。

Qwen3.5 与 MiMo 的 $D,F$ 互换，$D+F$ 与 $DF$ 均相同，
所以各档 $U$ 下的 $\qs$、$\qr$、RI 和末态有效容量相同。
它们的 gate/up 与 down 分项输入、以及共享扣除量不同；
汇总相同符合当前计数边界。

\par\smallskip\noindent\small
固定模型来源、精确字节及复算记录见本目录 README 与 data/；数学入口为 Table II(a)。
\end{document}
